\section{Natural higher brackets}
\label{sec.natural}

In order to prove the Theorem~\ref{thm.standardform} it is useful, and we believe interesting in itself, to consider the Koszul higher brackets
as special natural derived brackets: the notion of naturality has been discussed extensively in the paper \cite{markl1}; for our application we simply def\/ine a natural derived bracket of a linear endomorphism $f\colon A\to A$ as an element of type
\begin{gather*}\sum_{i+j=n}\psi_{ij} \mu_i\wedgebar (f\wedgebar\mu_j)\in D_n(A),\qquad \psi_{ij}\in \Q .\end{gather*}

We denote by
\begin{gather*}
\mathfrak{a}=\big\{\Psi\in \Hom_{\Q}(\Q[x],\Q[x])\,|\, \Psi(1)=0\big\}
=\left\{\sum_{n=0}^{\infty}p_n(x)\de^{n+1}\,|\, p_n(x)\in \Q[x]\right\}\end{gather*}
the associative algebra of linear endomorphisms of the $\Q$-vector space
$\Q[x]$ vanishing in~$1$, equipped with the composition product.
Every element $\Psi\in \mathfrak{a}$ is uniquely determined by its Taylor coef\/f\/icients,
i.e., by the inf\/inite double sequence $\psi_{ij}$, $i,j\ge 0$, def\/ined by the formulas:
\begin{gather}\label{equ.taylor}
\Psi(x^{i+1})=\sum_{j=0}^{\infty}\psi_{ij}\frac{x^j}{j!},\qquad i\ge 0,\qquad \psi_{ij}\in\Q .
\end{gather}

\begin{Definition} Let $\Psi\in \mathfrak{a}$ be a f\/ixed element.
For a commutative superalgebra $A$ over $\Q$ and a linear endomorphism $f\colon A\to A$,
the $\Psi$-bracket of $f$ is
\begin{gather*} \Psi_f\in D(A),\qquad
\Psi_f=\sum_{i,j\ge 0}\psi_{ij} \mu_j\wedgebar (f\wedgebar\mu_i),
\end{gather*}
where the $\psi_{ij}$'s are the Taylor coef\/f\/icients of $\Psi$ as in \eqref{equ.taylor}.
\end{Definition}

For every $n\ge i>0$, denote by $\Phi^{n,i}\in \mathfrak{a}$ the operator
\begin{gather*} \Phi^{n,i}(x^{i})=\frac{x^{n-i}}{(n-i)!} ,\qquad \Phi^{n,i}(x^s)=0\qquad \text{for}\quad s\not=i.\end{gather*}
Then for every $f\colon A\to A$ we have $\Phi^{n,i}_f=\mu_{n-i}\wedgebar (f\wedgebar\mu_{i-1})$:
\begin{gather*}\Phi^{n,i}_f(a_1,\ldots,a_{n})=\sum_{\sigma\in S(i,n-i)}\epsilon(\sigma)
f(a_{\sigma(1)}\cdots a_{\sigma(i)})a_{\sigma(i+1)}\cdots a_{\sigma(n)}.\end{gather*}
In particular, setting
\begin{gather*} \Phi^n=\sum_{i=1}^n(-1)^{n-i}\Phi^{n,i}\end{gather*}
we recover the def\/inition of higher Koszul brackets $\Phi^n_f$.
The formulas \eqref{equ.bracketrhoenne}, \eqref{equ.bracketmuenne} and \eqref{equ.bracketdenne} suggest the introduction of the Lie algebra
$\mathfrak{g}$ over the f\/ield $\Q$ of rational numbers with basis $l_0,l_1,l_2,\ldots$ and bracket
\begin{gather*}[l_n,l_m]=(n-m)\dfrac{(n+m+1)!}{(n+1)!(m+1)!}l_{n+m}.\end{gather*}
Notice that, setting $L_n=(n+1)! l_n$ we obtain the more familiar expression
\begin{gather*}[L_n,L_m]=(n-m)L_{n+m}.\end{gather*}
In particular, for every commutative superalgebra $A$ over $\Q$, the linear map
\begin{gather*} \mu_A\colon \ \mathfrak{g}\to D(A),\qquad \mu_A(l_n)=\mu_n\end{gather*}
is a morphism of Lie superalgebras: we denote by
\begin{gather*} \rho_A\colon \ \mathfrak{g}\to \Hom(D(A),D(A)),\qquad \rho_A(l_n)(\phi)=[\mu_n,\phi]\end{gather*}
the corresponding adjoint representation.

\begin{Lemma} The linear map $\rho\colon \mathfrak{g}\to \Hom(\mathfrak{a},\mathfrak{a})$:
\begin{gather*} \rho(l_n)(\Psi)=\left(\dfrac{x^n}{n!}-\dfrac{x^{n+1}\partial}{(n+1)!}\right)\circ \Psi-
\Psi\circ \dfrac{x\partial^{n+1}}{(n+1)!}\end{gather*}
is an injective morphism of Lie algebras.
\end{Lemma}

\begin{proof}We f\/irst prove that
\begin{gather*} \phi\colon \ \mathfrak{g}\to \Hom(\Q[x],\Q[x]),\qquad
\phi(l_n)= \dfrac{x\partial^{n+1}}{(n+1)!} ,\\
 \psi\colon \ \mathfrak{g}\to \Hom(\Q[x],\Q[x]),\qquad
\psi(l_n)=\dfrac{x^n}{n!}-\dfrac{x^{n+1}\partial}{(n+1)!} \end{gather*}
are morphisms of Lie algebras; it is easier to make the computation in the basis $L_n=(n+1)! l_n$,
\begin{gather*} \begin{split}
[\phi(L_n),\phi(L_m)]x^s&=\big[x\partial^{n+1},x\partial^{m+1}\big]x^s\\
&=x\partial^{n+1}\left(\prod_{i=0}^{m}(s-i)\right)x^{s-m}-x\partial^{m+1}\left(\prod_{i=0}^{n}(s-i)\right)x^{s-n}\\
&=(s-m)\prod_{i=0}^{m+n}(s-i)x^{s-m-n}-(s-n)\prod_{i=0}^{m+n}(s-i)x^{s-m-n}\\
&=(n-m)x\partial^{n+m+1}x^s=(n-m)\phi(L_{n+m})x^s.\end{split}\end{gather*}
Since $\psi(L_n)x^s=(n+1-s)x^{n+s}$ we have
\begin{gather*}\begin{split}
[\psi(L_n),\psi(L_m)]x^s&=\psi(L_n)(m+1-s)x^{m+s}-\psi(L_m)(n+1-s)x^{n+s}\\
&=((n+1-s-m)(m+1-s)-(m+1-n-s)(n+1-s))x^{n+m+s}\\
&=(n-m)(n+m+1-s)x^{n+m+s}=(n-m)\psi(L_{n+m})x^s.\end{split}\end{gather*}
Finally
\begin{align*}
[\rho(x),\rho(y)]\Psi&=\rho(x)\rho(y)\Psi-\rho(x)\rho(y)\Psi\\
&=\psi(x)\rho(y)\Psi-\rho(y)\Psi\phi(x)-
\psi(y)\rho(x)\Psi+\rho(x)\Psi\phi(y)\\
&=\psi(x) \psi(y)\Psi+\Psi\phi(y)\phi(x)
-\psi(y)\psi(x)\Psi-\Psi\phi(x)\phi(y)\\
&=[\psi(x),\psi(y)]\Psi-\Psi[\phi(x),\phi(y)]\\
&=\psi([x,y])\Psi-\Psi\phi([x,y])=\rho([x,y])\Psi .
\end{align*}
The injectivity is clear.
\end{proof}

\begin{Lemma} For every commutative superalgebra $A$ and every $f\in D_0(A)$ the map
\begin{gather*} \mathfrak{a}\to D(A),\qquad \Psi\mapsto \Psi_f\end{gather*}
is a morphism of $\mathfrak{g}$-modules. In other terms, for every $\Psi\in \mathfrak{a}$ and every $l\in \mathfrak{g}$ we have
\begin{gather}\label{equ.morfismogmoduli}
(\rho(l)\Psi)_f=\rho_A(l)(\Psi_f)=[\mu_A(l),\Psi_f] .
\end{gather}
\end{Lemma}

\begin{proof} By linearity is suf\/f\/icient to check \eqref{equ.morfismogmoduli} when $l=l_k$ and $\Psi=\Phi^{n,i}$. The proof follows from the following straightforward identities:
\begin{gather*}%\label{equ.actiononphienneni}
\rho(l_k)\Phi^{n,i}=\left(\binom{n-i+k}{k}-\binom{n-i+k}{k+1}\right)\Phi^{n+k,i}-\binom{k+i}{k+1}\Phi^{n+k,i+k},\\
\big[\mu_k,\Phi^{n,i}_f\big]=\left(\binom{n-i+k}{k}-\binom{n-i+k}{k+1}\right)\Phi^{n+k,i}_f-\binom{k+i}{k+1}\Phi^{n+k,i+k}_f.\tag*{\qed}
\end{gather*}
\renewcommand{\qed}{}
\end{proof}

\begin{Theorem}\label{thm.basisforPhienne}
In the notation above, for every integer $k>0$ let's denote by $\rho_k=\rho(l_k)\colon \mathfrak{a}\to \mathfrak{a}$.
For every integer $n>0$, there exists an unique sequence $c^n_1,\ldots,c^n_n$ of rational numbers such that
\begin{gather*}\Phi^{n+1}=\big(c^n_1\rho_1^n+c^n_2\rho_1^{n-2}\rho_2+c^n_3\rho_1^{n-3}\rho_3+\cdots+c^n_n\rho_n\big)\Phi^1.\end{gather*}
\end{Theorem}

\begin{proof}Denote by $V^n\subset \mathfrak{a}$ the subspace generated by $\Phi^{n,1},\ldots,\Phi^{n,n}$; notice that $\dim V^n=n$ and
$\rho_k(V^n)\subset V^{n+k}$. We claim that, for every $n\ge 2$ the $n$ vectors
\begin{gather*}\rho_1^{n-2}\big(\Phi^{2,2}\big),\qquad \rho_1^a\rho_{n-a}\big(\Phi^{1,1}\big),\qquad 0\le a\le n-2,\end{gather*}
form a basis of $V^{n}$. This is clear for $n=2$ since $\rho_1(\Phi^{1,1})=\Phi^{2,1}-\Phi^{2,2}$.
By induction on $n$, it is suf\/f\/icient to prove that for every $n\ge 2$ the linear map
\begin{gather*} V^{n}\oplus V^1\to V^{n+1},\qquad (x,y)\mapsto \rho_1(x)+\rho_n(y)\end{gather*}
is an isomorphism. Consider f\/irst the case $n=2$: in terms of matrix multiplication we have
\begin{gather*} \big(\rho_1\big(\Phi^{2,2}\big),\rho_1\big(\Phi^{2,1}\big),\rho_2\big(\Phi^{1,1}\big)\big)=\big(\Phi^{3,3},\Phi^{3,2},\Phi^{3,1}\big)
\begin{pmatrix}-3&0&-1\\
1&-1&0\\
0&1&1\end{pmatrix}\end{gather*}
and the determinant of the matrix is $\not=0$. Assume now $n\ge 3$: since
\begin{gather*} \rho_n\big(\Phi^{1,1}\big)=\Phi^{n+1,1}-\Phi^{n+1,n+1}
,\qquad\rho_1\big(\Phi^{n,n-2}\big)=-\binom{n-1}{2}\Phi^{n+1,n-1},\end{gather*}
we can write
\begin{gather*} \big(\rho_1\big(\Phi^{n,n}\big),\ldots,\rho_1\big(\Phi^{n,1}\big),\rho_n\big(\Phi^{1,1}\big)\big)=
\big(\Phi^{n+1,n+1},\ldots,\Phi^{n+1,1}\big)
\begin{pmatrix}A&B\\
0&C\end{pmatrix},\end{gather*}
where $A\in M_{3,3}(\Z)$ and $C\in M_{n-2,n-2}(\Z)$ are lower triangular matrices with nonzero entries in the diagonal, and therefore the determinant of the block matrix is $\not=0$. Thus, for every $n$ there exists an unique sequence of $n+1$ rational numbers $c^n_1,\ldots,c^n_n,b_n$ such that
\begin{gather*}\Phi^{n+1}=\big(c^n_1\rho_1^n+c^n_2\rho_1^{n-2}\rho_2+\cdots+c^n_n\rho_n\big)\Phi^1+
b_n\rho_1^{n-1}\big(\Phi^{2,2}\big)\end{gather*}
and we have to show that $b_n=0$. To this end it is suf\/f\/icient to consider the representation
\begin{gather*} \mathfrak{a}\to D(\Q[x]),\quad \Psi\mapsto \Psi_\de,\end{gather*}
where, as usual $\de=\frac{d~}{dt}\in D_0(\Q[x])$. Since $\de$ is a derivation we have
$\Phi^{n+1}_\de=[\mu_n,\de]=0$ for every $n\ge 0$, and therefore it is suf\/f\/icient to prove that $\rho_1^{n}(\Phi^{2,2}_\de)\not=0$ for every $n$.
By construction $\Phi^{2,2}_\de(t^a,t^b)=(a+b)t^{a+b-1}$ and an easy induction on $n$ shows that for every $n\ge 3$ we have
\begin{gather*} \rho_1^{n-2}\big(\Phi^{2,2}_\de\big)\big(t^{a_1},t^{a_2},\ldots,t^{a_n}\big)=
\left[\prod_{i=3}^n\binom{i-1}{2}\right](a_1+\cdots+a_n)t^{\sum a_i-1}.\tag*{\qed}\end{gather*}
\renewcommand{\qed}{}
\end{proof}
